\clearpage
\SourceLeaf{412}
\begin{center}\large ON THE TRANSFORMATION OF SURFACES.\end{center}
\noindent\hspace*{1.2em}The points of a surface given in itself, whose shape remains undetermined, are specified by $r$ and $\theta$. Here $r$ is the distance from a given point, measured along a shortest line, and $\theta$ is the angle that this shortest line makes with another fixed [shortest line] also passing through that point. Then every line for which $\theta$ is constant is cut at a right angle by every other line for which $r$ is constant. If an element of the latter line is put $=md\theta$, then $m$ depends only on the nature of the surface, independently of its shape, and is therefore a given function of $r$ and $\theta$.

Once a definite shape is assigned to the surface, the following holds for a line on which $\theta$ is constant.

Let $x,y,z$ be undetermined coordinates of a point on that line, and further let
\[
\begin{array}{lll}
\displaystyle\frac{\partial x}{\partial r}=\xi,&\displaystyle\frac{\partial y}{\partial r}=\eta,&\displaystyle\frac{\partial z}{\partial r}=\zeta,\\[6pt]
\displaystyle\frac{\partial\xi}{\partial r}=pX,&\displaystyle\frac{\partial\eta}{\partial r}=pY,&\displaystyle\frac{\partial\zeta}{\partial r}=pZ,
\end{array}
\]
\noindent where
\[
p=\BareRoot{\left(\frac{\partial\xi^2}{\partial r^2}+\frac{\partial\eta^2}{\partial r^2}+\frac{\partial\zeta^2}{\partial r^2}\right)},
\]
\noindent so that $\dfrac1p$ is the radius of curvature and
\[
XX+YY+ZZ=1.
\]
